\chapter{Conclusiones y trabajo futuro}
\label{chap:conclusiones}

\section{Objetivo 1: capacidad de ordenación}

La evidencia es favorable y robusta. El sistema alcanza Rank-IC medio \RankICSeleccion{} sobre
\CohortesSeleccion{} cohortes, con \FraccionICPositiva{} positivas, IC-IR \ICIRSeleccion{} y $t$ de
Newey--West \TNeweyWest{}. El test de permutación produce $p=\PValorPermutacion{}$ con
\NumPermutaciones{} réplicas; el bootstrap por bloques excluye cero; y la señal conserva el signo al
retirar eras completas, cambiar semillas, barajar etiquetas y neutralizar estilos conocidos. En la
era reservada el Rank-IC sigue siendo positivo (\RankICReservado{}).

La formulación defendible es que el sistema aprende una ordenación transversal fuera de muestra
dentro del universo, el periodo y el protocolo estudiados. No demuestra, en cambio, que la
arquitectura multiagente sea necesaria: el meta aprende a concentrar peso usando cohortes cerradas y
supera al promedio equiponderado, pero el agente de riesgo alcanza por sí solo un Rank-IC de
\RankICRisk{}, ligeramente superior. El resultado describe un meta que selecciona al mejor
especialista más que una combinación cuya superioridad haya quedado establecida.

La atribución local añade comprensión, no causalidad. Los huecos de apertura y la amplitud del rango
de negociación encabezan juntos \DosVariablesRisk{} de las observaciones de ese agente, y no son
simplemente beta ni volatilidad anual. El trabajo identifica así qué variables sostienen el
mecanismo, pero no puede distinguir si recogen liquidez, microestructura o una prima estable.

La lectura por eras también exige precisión. Que el Rank-IC sea mayor en el tramo donde el panel
histórico está más completo contradice una explicación simple y monotónica basada en la peor
cobertura temprana. No descarta que la disponibilidad del proveedor influya por otras vías. Por eso
la cobertura sigue siendo una amenaza alta y la réplica con otra fuente es trabajo futuro.

\section{Objetivo 2: cobrar el orden frente al mercado}

El segundo objetivo se juzga contra un listón relativo y de suma cero, no absoluto, y la respuesta
llega en dos mitades que conviene no fundir: las reglas de cartera sí alteran materialmente el
resultado, y la cartera elegida deja de perder frente al índice en la era reservada sin llegar a
superarlo con holgura.

\subsection*{2a: sensibilidad a la construcción de cartera}

Este subobjetivo queda respaldado por la rejilla completa. Manteniendo fijos modelo, señal, datos y
ventana, \NumCarteras{} combinaciones producen diferencias amplias de Information Ratio. El número
de posiciones separa las medianas en \InfluenciaPosiciones{}, la tenencia mínima en
\InfluenciaTenencia{} y el tope de efectivo en \InfluenciaEfectivo{}; la tolerancia de deriva y el
reparto de pesos mueven mucho menos el resultado.

La cartera elegida eleva el Information Ratio de \IRModelo{} a \IRCartera{} y el exceso geométrico
de \ExcesoModelo{} a \ExcesoCartera{}. Al mismo tiempo, reduce la rotación de \TurnoverModelo{} a
\TurnoverCartera{} veces al año y la máxima caída de \MDDModelo{} a \MDDCartera{}. No mejora por
reaccionar más, sino por respetar el horizonte de doce meses: amplía la tenencia mínima de medio
horizonte a uno completo y tolera más deriva antes de reequilibrar.

La magnitud de \IRCartera{} no debe aislarse de cómo fue obtenida. Es la mejor cifra de la rejilla y,
por tanto, una cota superior optimista. El hallazgo robusto del Objetivo 2a es la dispersión entre
políticas y el orden de importancia de las reglas, no el nivel exacto de la ganadora.

\subsection*{2b: comportamiento en la era reservada}

La comparación reservada es alentadora, pero preliminar. Con la cartera original, la señal produce
un exceso de \ExcesoModeloReservado{} e IR \IRModeloReservado{}. Con la política optimizada, sin
reentrenar ni cambiar el ranking, pasa a exceso \ExcesoCarteraReservado{} e IR
\IRCarteraReservado{}. El cambio de signo demuestra que la conclusión económica depende de la
cartera utilizada para medirla.

No demuestra generalización económica. La ventana contiene \CohortesReservadas{} cohortes cerradas
y \AnosCarteraReservada{} años de cartera; batir al índice en la mitad de los años significa uno de
dos. Además, la propia tenencia mínima reduce la rotación observable durante una serie tan corta.
El resultado es compatible con dejar de destruir valor y aproximadamente empatar al índice, no con
una ventaja sostenida.

La formulación final distingue así dos afirmaciones. Está demostrado que las reglas de cartera
alteran materialmente el resultado de selección. Está sugerido, pero no confirmado, que la política
elegida traduzca mejor la señal en una ventana nueva.

Sobre el objetivo tal como se enunció ---batir al índice--- la respuesta defendible es, por tanto,
negativa en su versión fuerte y afirmativa sólo en una versión débil: el sistema no demuestra
superar al mercado fuera de la ventana en la que se eligieron sus reglas, pero sí demuestra que
llegar a esa frontera dependía de la implementación tanto como del modelo. En un juego de suma cero,
pasar de destruir valor a empatar con el índice sin reentrenar nada es un resultado sobre dónde
estaba el cuello de botella, no una demostración de ventaja competitiva.

\section{Contribución defendible}

La contribución principal es un procedimiento auditable que separa información observable,
selección predictiva, pruebas de robustez e implementación económica. Esa disciplina permite
localizar una divergencia que una rentabilidad final habría ocultado: el orden transversal puede
mejorar mientras el pago de una cartera concreta empeora, y dos carteras pueden producir
conclusiones opuestas sobre la misma señal.

De ahí se sigue la tesis del trabajo: \textbf{ninguna afirmación práctica sobre un sistema
predictivo es interpretable sin declarar con qué cartera se midió}. La frase no convierte a la
cartera ganadora en recomendación de inversión. Obliga a tratar tamaño, efectivo, tenencia,
rebalanceo y costes como parte de la evaluación científica.

La evidencia económica tiene además una dependencia interna que las métricas agregadas no muestran:
la cartera pasó por \NumAccionesCartera{} acciones, pero una sola aporta \ContribucionAAPL{} puntos
porcentuales acumulados. Esa concentración reduce el número de experiencias independientes que
sostienen el resultado, y es la razón de que la conclusión económica deba ser más prudente que la
predictiva.

\section{Tres prioridades futuras}

\begin{enumerate}[leftmargin=*,itemsep=0.5em]
  \item \textbf{Acumular una era reservada con potencia real.} Es la única forma de comprobar si el
  cambio de signo económico persiste sin volver a utilizar 2015--2024. Más búsqueda sobre la misma
  ventana aumentaría la multiplicidad, no la evidencia independiente.
  \item \textbf{Replicar el protocolo con otra fuente y otro universo.} Una base histórica comercial
  completamente \textit{point-in-time} permitiría medir la sensibilidad a cobertura y comprobar si
  la ordenación se sostiene fuera de la gran capitalización estadounidense.
  \item \textbf{Modelar costes dependientes de liquidez y capacidad.} Una cartera de
  \PosicionesCartera{} posiciones requiere una función de impacto que complemente comisión y
  \textit{slippage} constantes y acote con precisión el patrimonio al que el resultado es aplicable.
\end{enumerate}

\section*{Respuesta final}

El TFM encuentra evidencia sólida de capacidad de ordenación y evidencia sólida de que la
construcción de cartera importa. No encuentra evidencia suficiente para afirmar que cinco agentes
sean necesarios ni que la cartera elegida generalice económicamente. Esa frontera entre lo
demostrado, lo sugerido y lo pendiente es el resultado que puede defenderse.
